\section{The proof-producing finite-trace compiler}
\label{sec:trace-compiler}

\subsection{Trace systems}

All run-time data are paired into one natural number \(u\): original inputs,
requested precision, numerical bounds, and externally supplied certificate
codes.  A \emph{trace system} consists of fixed primitive-recursive functions
\[
\begin{array}{lll}
 \mathsf{good}(u), & \mathsf{init}(u), & \mathsf{bound}(u),\\
 \mathsf{step}(i,s,u), & \mathsf{inv}(i,s,u), &
 \mathsf{accept}(s,u),
\end{array}
\]
where the Boolean-valued functions are normalized to \(0\) or \(1\).  Its run
is the primitive recursion
\begin{align*}
 \Run(0,u)&=\mathsf{init}(u),\\
 \Run(i+1,u)&=\mathsf{step}(i,\Run(i,u),u).
\end{align*}
In the repository's \(\PR\) syntax this is exactly
\(\code{PR.prec init step}\), with one parameter after the recursion argument.

The local proof obligations are open object-language formulas:
\begin{align}
 \mathsf{good}(u)=1 &\to
   \mathsf{inv}(0,\mathsf{init}(u),u)=1,
   \tag{B}\label{eq:trace-base}\\
 \mathsf{good}(u)=1\wedge\mathsf{inv}(i,s,u)=1 &\to
   \mathsf{inv}(i+1,\mathsf{step}(i,s,u),u)=1,
   \tag{S}\label{eq:trace-step}\\
 \mathsf{good}(u)=1\wedge
   \mathsf{inv}(\mathsf{bound}(u),s,u)=1 &\to
   \mathsf{accept}(s,u)=1.
   \tag{F}\label{eq:trace-final}
\end{align}
The essential requirement is not merely that these implications be true in
\(\N\).  A proof-producing trace package contains actual derivations
\[
 []\vdash[\PRAAx]\eqref{eq:trace-base},\qquad
 []\vdash[\PRAAx]\eqref{eq:trace-step},\qquad
 []\vdash[\PRAAx]\eqref{eq:trace-final},
\]
with the displayed free variables.  They may themselves be generated from
smaller proof-producing components.

\subsection{The compiler theorem}

\begin{theorem}[Embedded finite-trace reflection]
\label{thm:embedded-trace}
From object-language derivations of \eqref{eq:trace-base}--\eqref{eq:trace-final}
one can construct an object-language derivation
\[
 []\vdash[\PRAAx]
 \forall u\;\bigl(
  \mathsf{good}(u)=1\to
  \mathsf{accept}(\Run(\mathsf{bound}(u),u),u)=1
 \bigr).
\]
The construction is primitive recursive on the supplied proof terms.
\end{theorem}

\begin{proof}
We give the proof at the level of the repository's derivation constructors.
Keep \(u\) free and place \(G(u):\mathsf{good}(u)=1\) in the context.  Define
the induction formula with free variables \(i,u\):
\[
 A(i,u):\equiv
 \mathsf{inv}(i,\Run(i,u),u)=1.
\]
It is quantifier free because it is one equation between terms formed from
fixed primitive-recursive symbols.

For the base derivation, instantiate \eqref{eq:trace-base} at the current
\(u\), apply it to the context hypothesis \(G(u)\), and use the defining
zero equation
\[
 \Run(0,u)=\mathsf{init}(u)
\]
provided by \code{PRA.precZero}.  Equality substitution changes the resulting
invariant at \(\mathsf{init}(u)\) into \(A(0,u)\).

For the successor derivation, enter the induction context containing
\(A(i,u)\) and the shifted copy of \(G(u)\).  Instantiate
\eqref{eq:trace-step} with
\(s:=\Run(i,u)\).  Modus ponens gives
\[
 \mathsf{inv}(i+1,
   \mathsf{step}(i,\Run(i,u),u),u)=1.
\]
The successor defining equation supplied by \code{PRA.precSucc} is
\[
 \Run(i+1,u)=\mathsf{step}(i,\Run(i,u),u).
\]
Use \code{Prf.eqSubst1} (or simultaneous \code{Prf.eqSubst}) to rewrite the
state argument.  This produces \(A(i+1,u)\).

Apply \code{PRA.qfInd} to obtain, in context \(G(u)\),
\(\forall i\,A(i,u)\).  Eliminate the universal quantifier at the term
\(\mathsf{bound}(u)\).  Instantiate \eqref{eq:trace-final} with
\(s:=\Run(\mathsf{bound}(u),u)\) and apply it to \(G(u)\) and the specialized
invariant.  Discharge \(G(u)\) by implication introduction and close \(u\) by
universal introduction.

Every operation in this construction---substitution of terms into a supplied
derivation, weakening, implication elimination, equality substitution,
quantifier-free induction, and universal introduction---is a constructor or a
proved transformation in \code{ConciseIPRA.Deduction} and
\code{ConciseIPRA.PR}.  Consequently the construction can be implemented as a
Lean function producing a \(\Prf\) term, not merely as a metatheoretic
soundness argument.
\end{proof}

\begin{remark}[Why the theorem is nontrivial but modest]
The theorem does not assert that every true primitive-recursive statement is
provable in iPRA.  Its inputs already include proofs of the local arithmetic
obligations.  It states that bounded execution and invariant propagation add
no further logical strength: they are compiled by quantifier-free induction.
The mathematical burden is to design an invariant for which the local step is
actually derivable.
\end{remark}

\subsection{Finite query tables}

Suppose an input object is supplied through a fixed evaluator \(a(p,k)\) and a
global contract such as
\[
 H_a(p):\equiv\forall k\,\forall\ell\;
 R(a(p,k),a(p,\ell),k,\ell).
\]
A trace at requested precision \(n\) reads only finitely many samples.  Queries
may be adaptive: the next precision can depend on previous rational
comparisons.  Include in the trace state a query log and a Boolean accumulator
recording all instantiated contract facts.

Define \(\GoodTable(p,n,c)=1\) to be the finite conjunction of exactly the
regularity, domain, magnitude, and separation facts needed along the run coded
by \(c\).  If the algorithm computes its own query log, \(c\) may simply be the
run state.

\begin{lemma}[Finite instantiation]
\label{lem:finite-query}
For every fixed query generator, iPRA proves
\[
 H_a(p)\to \GoodTable(p,n,c)=1
\]
for the actual finite query table, provided each table field is an instance or
a finite algebraic consequence of \(H_a(p)\) and any other global hypotheses.
\end{lemma}

\begin{proof}
Induct over the finite query trace with a Boolean equation saying that all
entries recorded so far are valid.  At a successor stage, the next query
indices are primitive-recursive terms in the previous state.  Universal
elimination instantiates \(H_a(p)\) at those terms.  The resulting finite
arithmetic fact is inserted into the accumulator using a proved Boolean
conjunction lemma.  The induction formula contains only the accumulator
equation; the quantified hypothesis remains in the context.
\end{proof}

This is the exact repair for ``induction over all precisions.''  The requested
precision \(n\) is free, the run consumes a finite table, and only after the
fixed-\(n\) proof is complete is \(n\) universally introduced.

\begin{corollary}[Fix precision, then quantify]
\label{cor:fix-precision}
Suppose global hypotheses imply the finite input table for every free precision
\(n\), the trace laws are embedded derivations, and acceptance implies a finite
output contract \(\mathsf{Out}(p,n,A(p,n))\).  Then iPRA proves
\[
 H(p)\to\forall n\;\mathsf{Out}(p,n,A(p,n)).
\]
No induction formula contains \(\forall n\).
\end{corollary}

\subsection{Certificate mode}

A large computation need not be rerun by primitive recursion inside the final
theorem.  A certificate may contain the state sequence
\(s_0,\ldots,s_B\), local arithmetic witnesses, and the output.  A fixed
Boolean verifier checks
\begin{align*}
 s_0&=\mathsf{init}(u),\\
 s_{i+1}&=\mathsf{step}(i,s_i,u)\quad(i<B),\\
 \mathsf{inv}(i,s_i,u)&=1\quad(i\le B),\\
 \mathsf{accept}(s_B,u)&=1.
\end{align*}
Every clause is a bounded fold over the coded list.

\begin{theorem}[Certificate reflection]
\label{thm:certificate-reflection}
For a fixed certificate language and checker, one can construct an embedded
derivation
\[
 []\vdash[\PRAAx]\forall u\,\forall c\;
 \bigl(\TraceOK(u,c)=1\to
       \mathsf{Conclusion}(u,\mathsf{output}(c))\bigr),
\]
provided the local checker clauses have embedded soundness derivations.
\end{theorem}

\begin{proof}
Boolean-fold correctness exposes each checked list clause.  There are two
proof routes.  The direct route reads the final acceptance clause and applies
its embedded soundness theorem.  The reconstructive route proves by
quantifier-free induction on the list index that the certified state equals the
canonical run, then invokes \cref{thm:embedded-trace}.  The direct route is
usually smaller; the reconstructive route is useful when later theorems need
the canonical run equality.
\end{proof}

Certificate mode is appropriate for polynomial identities, isolating disks,
factorization traces, finite covers, high-precision analytic constants, and
large numerical inequalities.  A separate producer theorem may later show
that a certificate exists for a given input family.

\subsection{Sequential and network composition}

A proof-producing transducer \(F:X\to Y\) consists of a primitive-recursive
output algorithm, a finite input-demand function, and an embedded soundness
derivation.  If \(G:Y\to Z\) is another transducer, their composite first runs
\(F\), extracts the finite metadata needed to determine \(G\)'s demand, and
then runs \(G\).

\begin{theorem}[Composition]
\label{thm:trace-composition}
Proof-producing finite transducers are closed under sequential composition,
finite products, finite case splits with overlapping guards, and finite
acyclic networks.  The composite has primitive-recursive demand and bound
functions and an embedded soundness derivation.
\end{theorem}

\begin{proof}
For sequential composition, encode the state as a phase bit, a component state,
and frozen metadata.  Phase zero executes \(F\).  On its final step, the state
stores the certified output and all magnitude/domain data needed to initialize
\(G\).  Phase one executes \(G\).  The total bound is the sum of the two bounds
plus a switch step.

The invariant is a Boolean case split on the phase.  In phase zero it is
\(F\)'s invariant.  At the switch it includes \(F\)'s acceptance and \(G\)'s
input contract.  In phase one it is \(G\)'s invariant together with the frozen
metadata.  The local step proof is assembled from the component step proofs,
the switch soundness proof, and Boolean case analysis.  Apply
\cref{thm:embedded-trace}.  The final implication is \(G\)'s soundness after
substituting \(F\)'s certified output.

Products pair states and conjoin invariants.  Overlapping cases use a total
selector whose branch-correctness proof is supplied in each overlap.  A finite
acyclic network is topologically ordered and flattened into phases.  All proof
composition is substitution plus \code{Prf.trans} or ordinary implication
elimination.
\end{proof}

\subsection{The exact compiler boundary}

The trace compiler is a disciplined boundary between semantic mathematics and
formal arithmetic.

\begin{itemize}
\item Outside the boundary one proves local rational inequalities,
      coefficient identities, error budgets, and invariant preservation.
\item At the boundary each local theorem is expressed as a formula of
      \(\Fm\,\PTm\) and supplied as a \(\Prf\,\PRAAx\) derivation.
\item Inside the boundary generic code combines those derivations into the
      complete bounded-run theorem.
\item Only after that theorem exists is \code{pra\_iff\_native} applied.
\end{itemize}

This organization makes missing work visible.  A semantic Lean theorem can be
used as a guide for constructing a local derivation, but it cannot be inserted
as an unexamined premise.  Conversely, once the local obligations are embedded,
the generic control-flow proof should not be repeated theorem by theorem.
